\section{Implementation}

\subsection{Domain-Specific Constraint Language}

\paragraph{Constraint representation.}

The domain-specific constraint language was designed to represent the constraints of the original ILP formulation as relationships between computations and buffers.
Rather than encoding these relationships directly as integer variables and linear inequalities, the language gives each kind of constraint an explicit constructor.
The language is implemented by the algebraic data type \texttt{Constraint}.
The following code listing shows two examples of constraints that express ordering relationships between computations.

\begin{lstlisting}[language=Haskell]
data Constraint
  = ClusterBefore (Node Comp) (Node Comp)
  | ClusterBeforeUnlessFused (Node Comp) (Node Comp)
  -- Further constraints omitted...
\end{lstlisting}

The order of the computation clusters is represented by integer variables $\pi_i$, where $\pi_i$ denotes the cluster index of the $i$-th computation.
Constraints on these indices ensure that the cluster ordering respects the dependencies between computations, with fused computations sharing the same cluster index.

The constraint \texttt{ClusterBefore} expresses the strict ordering $\pi_i < \pi_j$ between two computations $i$ and $j$.
For dependencies that permit fusion, the constraint \texttt{ClusterBeforeUnlessFused} requires $\pi_i = \pi_j$ if the computations are fused, and $\pi_i < \pi_j$ otherwise.
This constraint makes the relationship between fusion decisions and cluster ordering explicit in the constraint language.

\paragraph{Lowering to ILP.}

To keep support for the external ILP solver, the constraints are lowered into a linear representation.
For each constraint, lowering produces a tuple containing linear constraints, variable bounds and a contribution to the objective.
The lowering signature and its result type are shown in the code listing below:

\begin{lstlisting}[language=Haskell]
lower :: LowerEnv -> Constraint -> Lower

type Lower = State String (LinearConstraint, Bounds, Expression)
\end{lstlisting}

The \texttt{LowerEnv} parameter provides the total number of computations in the problem, which is used for some constraints to generate appropriate bounds and coefficients during lowering.
The \texttt{Lower} type uses a state monad to supply fresh variable names to introduce auxiliary variables when needed.
Some constraints may not produce additional bounds or objective contributions.

For example, the \texttt{HorizontalReadCost} constraint requires auxiliary variables to encode the cost associated with distinct clusters reading the same buffer.
Fresh names ensure that variables introduced for separate constraints do not collide.

As a simple example of the translation, \texttt{ClusterBefore} is lowered to
\[
   \pi_i < \pi_j.
\]

The lowering of \texttt{ClusterBeforeUnlessFused} additionally uses a binary variable $x_{ij}$, where $x_{ij} = 0$ means that computations $i$ and $j$ are fused, and $x_{ij} = 1$ means they are not fused.
The constraint is encoded as:
\[
    x_{ij} \leq \pi_j - \pi_i \leq M \cdot x_{ij},
\]
where $M$ is a sufficiently large constant determined by the total number of computations in the problem.

The function \texttt{lowerAll} lowers all constraints using a shared supply of fresh variable names and combines the results.

\subsection{Presolver}

\subsection{Custom Solver}
